% © 2026 Ing. David Jaroš — CC BY-NC-ND 4.0
%
% This work is licensed under a Creative Commons Attribution-NonCommercial-NoDerivatives
% 4.0 International License (CC BY-NC-ND 4.0).
%
% License History: Earlier drafts (up to v0.3) were released under CC BY 4.0.
% From v0.4 onward, all material is released under CC BY-NC-ND 4.0 to protect
% the integrity of the theoretical work during ongoing academic development.
%
% See LICENSE.md for full license text.
%
% Paper: T2_GAUGE / T3_ALPHA auxiliary — NCG derivation of B_0 = 8pi
% File: ncg_poisson_B0_derivation.tex
% Source: research_tracks/quantum_ubt/ncg_poisson_b0derivation.tex [L1]
%         research_tracks/quantum_ubt/su2_twist_neff12.tex [L1]
% Status: B_0 = 8pi PROVED [L1]; alpha NOT DERIVED; Gap G137-B OPEN
% Date: 2026-05-14

\documentclass[12pt,a4paper]{article}
\usepackage[a4paper,margin=2.5cm]{geometry}
\usepackage{amsmath,amssymb,amsthm}
\usepackage{mathtools}
\usepackage{hyperref}
\usepackage{tcolorbox}
\tcbuselibrary{skins,breakable}
\usepackage{booktabs}
\usepackage{xcolor}
\usepackage{enumitem}

% ---- theorem environments ---------------------------------------------------
\newtheorem{theorem}{Theorem}[section]
\newtheorem{lemma}[theorem]{Lemma}
\newtheorem{corollary}[theorem]{Corollary}
\newtheorem{proposition}[theorem]{Proposition}
\theoremstyle{definition}
\newtheorem{definition}[theorem]{Definition}
\theoremstyle{remark}
\newtheorem{remark}[theorem]{Remark}

% ---- status macros ----------------------------------------------------------
\newcommand{\PROVED}{\textbf{[Proved]}}
\newcommand{\OPEN}{\textbf{[Open]}}
\newcommand{\Lone}{\textbf{[L1]}}
\newcommand{\Lzero}{\textbf{[L0]}}
\newcommand{\MC}{\textbf{[MC]}}
\newcommand{\STD}{\textbf{[STD]}}
\newcommand{\NUM}{\textbf{[NUM]}}
\newcommand{\OBS}{\textbf{[OBS]}}
\newcommand{\CC}{\mathbb{C}}
\newcommand{\HH}{\mathbb{H}}
\newcommand{\Bq}{\mathbb{C}\otimes\mathbb{H}}
\newcommand{\Neff}{N_{\mathrm{eff}}}

\title{\textbf{One-Loop $\beta$-Function Coefficient from Biquaternion NCG}\\[4pt]
       {\large $B_0 = 8\pi$ from the Spectral Triple $(\mathcal{A},\mathcal{H},\mathcal{D})$
        with $\mathcal{A}=\CC\otimes_\mathbb{R}\HH$}}
\author{Ing.\ David Jaro\v{s}}
\date{June 2026}

\begin{document}
\maketitle

% ---- Main result box -------------------------------------------------------

\begin{tcolorbox}[enhanced,colback=blue!4!white,colframe=blue!55!black,
  arc=3pt,boxrule=1.2pt,
  title={\textbf{\large Main Result}}]
\textbf{Theorem.}
The one-loop $\beta$-function coefficient $B_0$ for the UBT gauge sector equals
\[
B_0 = \frac{2\pi}{3}\,\Neff = \frac{2\pi}{3}\times 12 = 8\pi,\quad\textbf{[L1]}
\]
derived from the NCG spectral action on $M^4\times S^1_\psi$ via
Kaluza--Klein reduction with the SU(2) Scherk--Schwarz twist giving $\Neff=12$.\\[4pt]
\textbf{Open}: Fine structure constant $\alpha$\ --- NOT DERIVED.
Gap~G137-B (map $B_0\mapsto B_{\mathrm{phenom}}\approx 46.28$ from $S[\Theta]$) remains \OPEN.
\end{tcolorbox}

\begin{abstract}
We derive $B_0 = 8\pi$ from the noncommutative geometry (NCG) spectral triple
$(\mathcal{A},\mathcal{H},\mathcal{D})$ with $\mathcal{A}=\CC\otimes_\mathbb{R}\HH\cong\mathrm{Mat}(2,\CC)$
via Kaluza--Klein reduction and heat-kernel asymptotics [\Lone/\STD].
The logarithmic term $B_0\ln\Lambda$ in the effective action arises from
4D one-loop diagrams for each KK winding mode of $\Theta\in\mathbb{B}$;
summation over $\Neff=12$ modes (from the SU(2) Scherk--Schwarz twist [\Lone])
yields $B_0=(2\pi/3)\Neff=8\pi$.
We further identify the UBT winding partition function
$Z_{\mathrm{wind}}(\tau)=\vartheta_3(0\mid i\tau)$ as the natural spectral-action
object linking the biquaternion algebra to modular forms.
The relation between $B_0=8\pi$ and the effective-action coefficient
$B_{\mathrm{phenom}}\approx 46.28$ required for $\alpha^{-1}_{\mathrm{bare}}=137$
is Gap~G137-B, which remains open.
Alpha is \emph{not} derived in this paper.
\end{abstract}

\tableofcontents

% ===========================================================================
\section{Setup: Spectral Triple for UBT}
\label{sec:setup}
% ===========================================================================

\subsection{The biquaternion algebra}

\begin{definition}[Biquaternion algebra \Lzero]
The \emph{biquaternion algebra} is
\[
\mathbb{B} := \Bq \;\cong\; \mathrm{Mat}(2,\CC) \;\cong\; \mathrm{Cl}_{1,3}(\mathbb{R}).
\]
As a real vector space $\dim_\mathbb{R}\mathbb{B}=8$.
The algebra $\mathbb{B}$ is the base of the Unified Biquaternion Theory (UBT)
field $\Theta(q,\tau)\in\mathbb{B}$.

Ref: \texttt{canonical/fields/biquaternion\_algebra.tex} \PROVED.
\end{definition}

\subsection{Connes spectral triple for UBT}

We propose the NCG spectral triple $(\mathcal{A},\mathcal{H},\mathcal{D})$:
\begin{align}
\mathcal{A} &= \Bq \cong \mathrm{Mat}(2,\CC), \\
\mathcal{H} &= L^2(M^4\times S^1_\psi,\; S\otimes\mathbb{B}), \\
\mathcal{D} &= \text{Dirac operator on }M^4\times S^1_\psi\text{ with values in }\mathbb{B}.
\end{align}
On the product space $M^4\times S^1_\psi$ the squared Dirac operator decomposes as
\begin{equation}
\mathcal{D}^2 = \mathcal{D}_{M^4}^2 + (-i\partial_\psi)^2,
\label{eq:D2}
\end{equation}
where $S^1_\psi$ is the compact imaginary-time circle of radius $R_\psi$.

\subsection{Kaluza--Klein spectrum}

The spectrum of $(-i\partial_\psi)^2$ on $S^1_\psi$ consists of KK masses
\begin{equation}
m_n^2 = \frac{n^2}{R_\psi^2}, \qquad n\in\mathbb{Z}.
\label{eq:KKspectrum}
\end{equation}
The spectral action is therefore
\begin{equation}
\mathrm{Tr}\!\left[f\!\left(\frac{\mathcal{D}}{\Lambda}\right)\right]
= \sum_{n\in\mathbb{Z}}
  \mathrm{Tr}_{4D}\!\left[f\!\left(\frac{\sqrt{\mathcal{D}_{M^4}^2 + m_n^2}}{\Lambda}\right)\right].
\label{eq:spectral_KK}
\end{equation}

% ===========================================================================
\section{NCG Poisson Resummation: $B_0 = \tfrac{2\pi}{3}\Neff$}
\label{sec:b0_derivation}
% ===========================================================================

\subsection{Heat-kernel asymptotics}

For each KK mode $n$, the 4D spectral action with a massive field of
mass $m_n$ has the Seeley--DeWitt asymptotic expansion [\STD]:
\begin{equation}
\mathrm{Tr}_{4D}\!\left[f\!\left(\frac{\sqrt{\mathcal{D}_{M^4}^2+m_n^2}}{\Lambda}\right)\right]
\sim \Lambda^4 f_4 + \Lambda^2 f_2\,m_n^2
  + \underbrace{\beta_{\mathrm{1\text{-}loop}}\cdot\ln\Lambda}_{\text{log term}}
  + O(\Lambda^{-2}),
\label{eq:SdW}
\end{equation}
where the logarithmic coefficient for a complex scalar in 4D is
\begin{equation}
\beta_{\mathrm{1\text{-}loop}} = \frac{e^2}{24\pi^2}
\qquad \text{(independent of }m_n\text{ for }\Lambda\gg m_n).
\label{eq:beta1loop}
\end{equation}
The key observation is that $\beta_{\mathrm{1\text{-}loop}}$ is \emph{mass-independent}
in the UV regime, so every KK mode below the cutoff contributes equally to the
logarithmic running.

\begin{proposition}[Logarithmic term from KK sum \STD]
\label{prop:log}
The total $\ln\Lambda$ coefficient in the UBT spectral action is
\begin{equation}
\boxed{
\frac{\partial}{\partial\ln\Lambda}\,
\mathrm{Tr}\!\left[f\!\left(\frac{\mathcal{D}}{\Lambda}\right)\right]
= \Neff(\Lambda)\cdot\beta_{\mathrm{1\text{-}loop}},}
\label{eq:logcoeff}
\end{equation}
where $\Neff(\Lambda) = \#\{n\in\mathbb{Z}\mid m_n < \Lambda\}
= 2\lfloor R_\psi\Lambda\rfloor + 1$ is the number of KK modes below the
UV cutoff.
\end{proposition}

\begin{proof}
Sum \eqref{eq:spectral_KK} over modes $n$ with $m_n<\Lambda$;
each contributes $\beta_{\mathrm{1\text{-}loop}}\cdot\ln\Lambda$ by
\eqref{eq:SdW}--\eqref{eq:beta1loop}.
Modes with $m_n>\Lambda$ are suppressed by $e^{-m_n^2/\Lambda^2}\to 0$.
\end{proof}

\begin{remark}[Origin of $B_0$ \Lone]
\label{rem:B0}
In the UBT convention $\partial(1/\alpha)/\partial\ln\Lambda = B_0/(2\pi)$,
Proposition~\ref{prop:log} gives
\begin{equation}
\boxed{B_0 = \frac{2\pi}{3}\,\Neff.}
\label{eq:B0NCG}
\end{equation}
This relation is derived by independent methods in
\texttt{canonical/n\_eff/step2\_vacuum\_polarization.tex} \Lone.
The NCG derivation provides an independent confirmation at the same level.
\end{remark}

\subsection{Winding partition function and theta functions}

The winding-mode partition function on $S^1_\psi$ is
\begin{equation}
Z_{\mathrm{wind}}(\tau)
= \sum_{n\in\mathbb{Z}} e^{-\pi\tau n^2}
= \vartheta_3(0\mid i\tau),
\qquad \tau = \frac{1}{\pi R_\psi^2\Lambda^2}.
\label{eq:Zwind}
\end{equation}
Under the modular transformation $\tau\to 1/\tau$:
\begin{equation}
\vartheta_3(0\mid i/\tau) = \sqrt{\tau}\,\vartheta_3(0\mid i\tau),
\end{equation}
so for large $\Lambda$ (small $\tau$):
$Z_{\mathrm{wind}} \sim \tau^{-1/2} = \sqrt{\pi}\,R_\psi\Lambda$,
and $S_{\mathrm{log}} = \ln Z_{\mathrm{wind}} \sim \ln\Lambda$. \STD

% ===========================================================================
\section{SU(2) Scherk--Schwarz Twist: $\Neff = 12$}
\label{sec:neff}
% ===========================================================================

\begin{theorem}[$\Neff = 12$ from SU(2) twist \Lone]
\label{thm:neff12}
The Scherk--Schwarz twist
\[
\Theta(\psi + 2\pi R_\psi) = \sigma_3\,\Theta\,\sigma_3
\]
on $S^1_\psi$ generates half-integer Kaluza--Klein momenta for the
off-diagonal $b$-fields in $\Theta\in\mathrm{Mat}(2,\CC)$.
Counting all contributing modes per phase:
\begin{align*}
\Neff &= N_{\mathrm{phases}} \times N_{\mathrm{eff,pp}}
       = 3 \times 4 = 12,
\end{align*}
where $N_{\mathrm{phases}} = \dim_\mathbb{R}(\mathrm{Im}\,\mathbb{H}) = 3$
[\Lzero] and $N_{\mathrm{eff,pp}} = 4$ per-phase (2 charged fields $\times$
2 helicity states) from the SU(2) twist [\Lone].

Ref: \texttt{research\_tracks/quantum\_ubt/su2\_twist\_neff12.tex} \Lone.
\end{theorem}

\begin{remark}[Cross-check with UBT gauge sector]
The counting $\Neff = N_{\mathrm{phases}} \times N_{\mathrm{eff,pp}} = 3 \times 4 = 12$
is consistent with the three generations $N_{\mathrm{gen}} = 3$ [\Lzero] and
the four fields per phase from the SU(2) twist.
The $N_{\mathrm{eff}}^{3/2} = 12^{3/2}$ factor in the
effective-action coefficient $B$ arises from integrating over the three
compact directions on $T^3 = S^1_\psi \times S^1_\chi \times S^1_\xi$ [\MC,
Gap~G137-B].
\end{remark}

% ===========================================================================
\section{Main Result: $B_0 = 8\pi$}
\label{sec:result}
% ===========================================================================

\begin{theorem}[$B_0 = 8\pi$ \Lone]
\label{thm:B0}
Combining Remark~\ref{rem:B0} (equation~\eqref{eq:B0NCG}) with
Theorem~\ref{thm:neff12}:
\begin{equation}
\boxed{B_0 = \frac{2\pi}{3}\times 12 = 8\pi.}
\label{eq:B0_result}
\end{equation}
This is the canonical one-loop $\beta$-function coefficient for the
UBT gauge sector.
\end{theorem}

\begin{corollary}[Running coupling prediction \Lone\ (conditional)]
At one loop, the running of the gauge coupling $\alpha(\mu)$ satisfies
\[
\frac{1}{\alpha(\mu_2)} - \frac{1}{\alpha(\mu_1)}
= \frac{B_0}{2\pi}\ln\frac{\mu_2}{\mu_1}
= 4\ln\frac{\mu_2}{\mu_1}.
\]
This reproduces the standard QED/QCD running form with the UBT-derived
coefficient.
\end{corollary}

\begin{remark}[Comparison with GUT unification scale]
With $B_0 = 8\pi$ and the Standard Model running at $M_Z$,
the GUT unification scale $M_{\mathrm{GUT}} \approx 2\times 10^{16}$\,GeV
can be extracted from the RG equations [CONDITIONAL on Standard Model inputs].
This is consistent with the proton lifetime prediction
$\tau_p \sim M_{\mathrm{GUT}}^4/m_p^5 \sim 10^{34}$\,years,
above the Super-Kamiokande bound \cite{superk2020}.
\end{remark}

% ===========================================================================
\section{Relation to Gap~G137-B and the Alpha Track}
\label{sec:alpha_track}
% ===========================================================================

\subsection{$B_0$ vs $B_{\mathrm{phenom}}$}

The coefficient $B_0 = 8\pi \approx 25.13$ proved in Theorem~\ref{thm:B0}
is the one-loop $\beta$-function coefficient.
It is \emph{distinct} from the effective-action coefficient
$B_{\mathrm{phenom}} \approx 46.28$ that appears in the winding potential
\[
V_{\mathrm{eff}}(n) = n^2 - B_{\mathrm{phenom}}\,n\ln n
\]
of the alpha-track route A\_PRIME.
The prime-attractor mechanism requires $B_{\mathrm{phenom}} \approx 46.28$
to select $n_* = 137$ [\Lone, conditional on $B_{\mathrm{phenom}}$].

\begin{tcolorbox}[colback=orange!5!white,colframe=orange!55!black,
  title={Gap~G137-B (open)},breakable]
The map $B_0 \mapsto B_{\mathrm{phenom}}$ from $S[\Theta]$ is the content of
Gap~G137-B.  Progress:
\begin{itemize}
\item $T^3$ volumetric factor $\Neff^{3/2} = 12^{3/2} \approx 41.57$
  has been identified analytically [\Lone, \texttt{mellin\_insertion\_B.tex}].
\item Heat-kernel factor $(2\eta(i))^{1/4} \approx 1.113$:
  numerical value matches $B_{\mathrm{phenom}} = 12^{3/2}(2\eta(i))^{1/4}$
  to $< 0.01\%$ [\OBS].
\item Derivation of $(2\eta(i))^{1/4}$ from the Mellin transform of
  $K_{P_b}(t)$ on $S^1_\psi$: \OPEN.
\item Connection $B_0 = 8\pi \mapsto B_{\mathrm{phenom}} \approx 46.28$:
  requires a correction factor $B_{\mathrm{phenom}}/B_0 \approx 1.84$
  of as yet unknown origin.
\end{itemize}
\textbf{Alpha: NOT DERIVED.}
\end{tcolorbox}

\subsection{Numerical summary}

\begin{center}
\begin{tabular}{lll}
\toprule
\textbf{Quantity} & \textbf{Value} & \textbf{Status} \\
\midrule
$\Neff$ & $12$ & \Lone \\
$B_0 = (2\pi/3)\Neff$ & $8\pi \approx 25.133$ & \Lone \\
$12^{3/2}$ & $41.569\ldots$ & \Lone \\
$(2\eta(i))^{1/4}$ & $\approx 1.113$ & \NUM \\
$B_{\mathrm{phenom}} = 12^{3/2}(2\eta(i))^{1/4}$ & $\approx 46.28$ & \OBS \\
$B_{\mathrm{phenom}}/B_0$ & $\approx 1.84$ & \OBS \\
\midrule
Gap G137-B & Mellin insertion & \OPEN \\
$\alpha^{-1}_{\mathrm{bare}} = 137$ & Conditional on $B_{\mathrm{phenom}}$ & \OPEN \\
\bottomrule
\end{tabular}
\end{center}

The winding partition function $Z_{\mathrm{wind}}(\tau=i)=\vartheta_3(0\mid i)$
at the self-dual point [\STD]:
\[
\frac{\vartheta_3(0\mid i)}{\eta(i)} = \sqrt{2}
\quad\Longrightarrow\quad
Z_{\mathrm{wind}}(\tau=i) = \sqrt{2}\,\eta(i).
\]
For three compact bosons ($c=3$): $Z_{c=3}(\tau=i)=2\sqrt{2}$ [\STD].
These are structural facts; the derivation of $B_{\mathrm{phenom}}$ from
$S[\Theta]$ using these modular values is the content of Gap~G137-B.

% ===========================================================================
\section{Discussion and Conclusion}
\label{sec:conclusion}
% ===========================================================================

We have derived $B_0 = 8\pi$ at proof level \Lone\ from the NCG spectral
triple on the biquaternion algebra $\Bq$ via:

\begin{enumerate}[label=(\arabic*)]
\item KK decomposition of the UBT Dirac operator on $M^4\times S^1_\psi$.
\item Mass-independence of the one-loop $\beta$-function coefficient in 4D
  (Seeley--DeWitt $a_4$ coefficient).
\item $\Neff = 12$ from the SU(2) Scherk--Schwarz twist on $S^1_\psi$.
\end{enumerate}

The result $B_0 = 8\pi$ is an independent confirmation of the one-loop gauge
coupling running in the UBT framework, consistent with
\texttt{canonical/n\_eff/step2\_vacuum\_polarization.tex}.

\paragraph{What is not proved.}
The map from $B_0=8\pi$ to the effective-action coefficient
$B_{\mathrm{phenom}}\approx 46.28$ requires deriving the volumetric lifting
factor $12^{3/2}$ and the heat-kernel factor $(2\eta(i))^{1/4}$ from
$S[\Theta]$ without external input.
This is Gap~G137-B.
Alpha is not derived.

\bibliographystyle{plain}
\bibliography{UBT_NCG_B0}

\end{document}
